\section{Limitations and responsible use}

The evidence base is purposive, single-coder, and limited to one public specialist community. It was constructed to maximize conceptual coverage and metric design, not representativeness. Public discussion is shaped by installed base, product complexity, community participation, shareability, and migration to private support. Open thread status cannot be equated with operational non-resolution.

The field-level case sets are small and heterogeneous. Their scores demonstrate operationalization and expose missing fields; they are not calibrated probabilities or common performance measures. Pairwise associations arise from multi-label coding in a purposive corpus and remain descriptive. Source-reported product evidence has not been converted into independent replication. The community schemas are designed artifacts whose operational value remains to be evaluated.

Responsible reuse requires preservation of the original analytical unit, denominator, uncertainty state, validation stage, and prohibited inference. Users should retain source provenance and contributor ownership, avoid redistributing copied discussions, and document corrections through explicit versioning and lineage. Extension to another domain should include domain-specific code boundaries and independent validation instead of assuming that the ten constructs are exhaustive.
